% Part: incompleteness
% Chapter: introduction
% Section: overview

\documentclass[../../../include/open-logic-section]{subfiles}

\begin{document}

\olfileid[pt]{inc}{int}{ovr}

\olsection{Panorama dos resultados de incompletude}

Hilbert esperava que a matemática pudesse ser formalizada em uma
teoria !!{axiomatizable} que fosse possível provar !!{complete} e
!!{decidable}. Além disso, pretendia provar a consistência dessa teoria
por meios muito fracos, ``finitários'', o que defenderia a matemática
clássica dos desafios do intuicionismo. Os teoremas da incompletude de
G\"odel mostraram que esses objetivos não podem ser alcançados.

O primeiro teorema da incompletude de G\"odel mostrou que uma versão
dos \emph{Principia Mathematica} de Russell e Whitehead não é
!!{complete}. Mas a prova era, na verdade, muito geral e se aplica a
uma ampla variedade de teorias. Isso significa que não foi apenas o
caso de os \emph{Principia Mathematica} não conseguirem captar
completamente a matemática: \emph{nenhuma} teoria aceitável consegue.
Levou algum tempo para isolar as características das teorias que bastam
para a aplicação dos teoremas da incompletude e generalizar a prova de
G\"odel de modo que dependesse apenas delas. Agora, porém, estamos em
condições de enunciar uma versão muito geral do primeiro teorema da
incompletude para teorias na linguagem aritmética $\Lang L_A$.

\begin{thm}
Se $\Gamma$ é uma teoria consistente e !!{axiomatizable} em~$\Lang
L_A$ que !!{represents} todas as funções computáveis e relações
!!{decidable}, então $\Gamma$ não é !!{complete}.
\end{thm}

Dizer que $\Gamma$ não é !!{complete} é dizer que, para pelo menos uma
!!{sentence}~$!A$, $\Gamma \Proves/ !A$ e $\Gamma \Proves/ \lnot !A$.
Tal !!a{sentence} é chamada de \emph{independente} (de~$\Gamma)$. Na
verdade, podemos provar com relativa rapidez que devem existir
sentenças independentes. Mas a força da prova do teorema de G\"odel
reside no fato de que ela apresenta um \emph{exemplo específico} de
tal !!{sentence} independente. A engenhosa construção produz
!!a{sentence}~$!G_\Gamma$, chamada de \emph{sentença de G\"odel}
para~$\Gamma$, que não é demonstrável porque, em $\Gamma$,
$!G_\Gamma$ equivale à afirmação de que $!G_\Gamma$ não é demonstrável
em~$\Gamma$. A construção faz isso \emph{construtivamente}, isto é,
dadas uma axiomatização de~$\Gamma$ e uma descrição do sistema de
!!{derivation}, a prova fornece um método para escrever efetivamente
$!G_\Gamma$.

A construção da prova de G\"odel exige que encontremos uma maneira de
expressar em $\Lang L_A$ as propriedades de termos e !!{formula}s da
própria $\Lang L_A$, bem como as operações sobre eles. Entre essas
propriedades estão ``$!A$ é !!a{sentence}'' e ``$\delta$ é
!!a{derivation} de~$!A$'', e entre as operações está
$\Subst{!A}{t}{x}$. Essa maneira deve (a) expressar tais propriedades e
relações mediante uma ``codificação'' dos símbolos e de suas sequências
(que é o que são termos, !!{formula}s, !!{derivation}s etc.) como
números naturais (aquilo de que $\Lang L_A$ pode falar). Deve também
(b) fazer isso de tal modo que $\Gamma$ prove os fatos pertinentes;
portanto, precisamos mostrar que essas propriedades são codificadas por
propriedades !!{decidable} dos números naturais e que as operações
correspondem a funções computáveis sobre os números naturais. Isso se
chama ``aritmetização da sintaxe''.

Antes de investigar como a sintaxe pode ser aritmetizada, porém,
consideraremos a condição de que~$\Gamma$ seja ``forte o suficiente'',
isto é, de que !!{represents} todas as funções computáveis e relações
!!{decidable}. Isso exige uma definição precisa de ``computável''. Ela
pode ser dada de várias maneiras, por exemplo, pelo modelo das máquinas
de Turing ou como as funções computáveis por programas em alguma
linguagem de programação de uso geral. Contudo, como nosso objetivo é
!!{represents} essas funções e relações em uma teoria na
linguagem~$\Lang L_A$, convém escolher uma definição simples de
computabilidade apenas para funções numéricas. Trata-se da noção de
\emph{função recursiva}. Assim, discutiremos primeiro as funções
recursivas. Em seguida, mostraremos que $\Th{Q}$ já !!{represents}
todas as funções e relações recursivas. Isso nos permitirá aplicar o
teorema da incompletude a teorias específicas, como $\Th{Q}$ e
$\Th{PA}$, pois teremos estabelecido que elas são exemplos de teorias
``fortes o suficiente''.

O resultado final da aritmetização da sintaxe é !!a{formula}
$\OProv[\Gamma](x)$ que, mediante a codificação de !!{formula}s como
números, expressa a demonstrabilidade a partir dos axiomas
de~$\Gamma$. Mais especificamente, se $!A$ é codificada pelo
número~$n$ e $\Gamma \Proves !A$, então
$\Gamma \Proves \Prov[\Gamma](\num{n})$. Esse ``predicado de
demonstrabilidade'' de $\Gamma$ também nos permite expressar, em certo
sentido, a consistência de $\Gamma$ como !!a{sentence} de~$\Lang L_A$:
seja a ``afirmação de consistência'' de~$\Gamma$ a !!{sentence}
$\lnot \Prov[\Gamma](\num{n})$, na qual tomamos $n$ como o código de
uma contradição, por exemplo, de~$\lfalse$. O segundo teorema da
incompletude afirma que teorias consistentes e !!{axiomatizable}
também não provam suas próprias afirmações de consistência. As
condições necessárias para que esse teorema se aplique são um pouco
mais exigentes do que simplesmente a teoria representar todas as
funções computáveis e relações !!{decidable}, mas mostraremos que
$\Th{PA}$ as satisfaz.

\end{document}
